\documentclass[12pt, letterpaper]{article}

\usepackage[margin=1in]{geometry}
\usepackage{graphicx}
\usepackage{float}
\usepackage{booktabs}
\usepackage{parskip}
\usepackage{setspace}
\usepackage{caption}
\usepackage{subcaption}
\usepackage{hyperref}
\usepackage{xcolor}
\usepackage{amsmath}

\graphicspath{
  {figures/}
  {figures/explore/}
  {figures/explore/apc/}
  {figures/explore/matrices/}
  {figures/sex/}
  {figures/political/}
  {figures/nonblack/}
  {figures/bp/}
  {figures/gss-decade/}
  {figures/nativity/}
  {figures/homogeneity/}
  {figures/catholic/}
  {figures/demographics/}
  {figures/validation/}
}

\captionsetup{font=small, labelfont=bf, skip=4pt}
\onehalfspacing

\newcommand{\todo}[1]{\textcolor{red}{\textbf{[TODO: #1]}}}
\newcommand{\placeholder}[1]{\textcolor{blue}{\textbf{[PLACEHOLDER: #1]}}}

\title{Religious Persistence in the United States\\
  \large Results Draft}
\author{Renato de Angelis}
\date{\today}

%% ─────────────────────────────────────────────────────────────────────────────

\begin{document}
\maketitle
\tableofcontents
\clearpage

Two headline patterns characterize the intergenerational transmission of religious
affiliation across birth cohorts from 1925 to 1984. First, aggregate persistence
has declined: the second-largest eigenvalue (in modulus) of the cohort-specific
transition matrix ($\lambda_2$) falls approximately linearly from 0.80 in the
earliest cohort window to 0.60 in the latest---a result robust to age
standardization and to the choice of cohort binning. Second, this decline reflects
a structural reorientation of American religious life toward two poles. The
religiously unaffiliated (``None'') and Evangelical Protestants are simultaneously
the most stable destinations for exiters from other traditions and, increasingly,
stable identities in their own right. Catholic and Mainline Protestant traditions,
by contrast, are losing retention without gaining converts: both receive negligible
inflow in every cohort examined, and their implied long-run shares contract
accordingly.

The demographic heterogeneity results indicate that exit into nonaffiliation is
more pronounced among men and political liberals, and that recent cohorts show a
degree of religious-political sorting across all strata. These differences are
substantively interesting but do not alter the structural conclusion: across
subgroups, the basic reorientation toward a None--Evangelical bipolarity is
consistent.

%% ═══════════════════════════════════════════════════════════════════════════
\section{National Trend}
%% ═══════════════════════════════════════════════════════════════════════════

\subsection{Transition matrices by birth cohort}

\begin{figure}[H]
  \centering
  \includegraphics[width=\textwidth]{matrices/P_grid_national_10yr.png}
  \caption{RELIG16 $\to$ RELIG transition matrices by 10-year birth cohort,
           1925--1984. Rows = origin religion (RELIG16); columns = current
           religion (RELIG). Cell values are row-conditional probabilities.}
  \label{fig:P_grid}
\end{figure}

\subsection{Origin and implied stationary distributions}

\begin{figure}[H]
  \centering
  \includegraphics[width=\textwidth]{matrices/pi_dist_grid_10yr.png}
  \caption{Share of each religious group in the origin distribution ($\pi_0$,
           top row) and the implied stationary distribution ($\pi^*$, bottom
           row), by 10-year birth cohort. $\pi^*$ is the long-run distribution
           implied by cohort $t$'s regime if that regime were to persist.}
  \label{fig:pi_dist}
\end{figure}

The cohort-specific transition matrices (Figure~\ref{fig:P_grid}) display the
cell-level shifts that produce the aggregate trend. Reading across cohorts, the
dominant change is the growth of the ``None'' column: outflows from all affiliated
origin states into nonaffiliation rise substantially in each successive 10-year
window, with the ordering Other $>$ Mainline $>$ Catholic $>$ Evangelical
preserved throughout. Mid-single-digit outflows in the 1925 cohort expand to
rates in the twenties and thirties by the 1975 cohort. Against this common trend,
origin states diverge sharply in their diagonals. Catholic persistence falls from
approximately 0.85 to 0.62 across the full range; Evangelical persistence is
notably stable (0.74--0.80). Conversion into Catholic is negligible in all
cohorts; by the post-1950 cohorts, Evangelical is effectively the only non-None
destination receiving meaningful inflow from other traditions. Outflow from None
into Evangelical remains relatively steady even as None's own diagonal persistence
rises dramatically---a pattern suggesting that a stable subpopulation continues to
move between these two poles while both categories simultaneously strengthen.
Figure~\ref{fig:pi_dist} traces the implied stationary distributions ($\pi^*$):
moving from the 1925 to the 1975 cohort, $\pi^*$ shifts substantially toward None
and away from Catholic and Mainline, while Evangelical holds relatively stable.
These are conditional summaries---what each cohort's regime implies for the long
run if that regime were to persist---not forecasts. The widening gap between
$\pi_0$ and $\pi^*$ in later cohorts indicates that each successive regime implies
a more dramatic departure from its own starting composition.

\clearpage

%% ═══════════════════════════════════════════════════════════════════════════
\section{Memory Measures}
%% ═══════════════════════════════════════════════════════════════════════════

\subsection{Individual memory curves}

\begin{figure}[H]
  \centering
  \includegraphics[width=0.95\textwidth]{memory/im_memory_10yr.png}
  \caption{Individual memory (log TV distance from $\pi^*$) by origin state
           and generation step $t = 0$--$4$, for each 10-year birth cohort.
           Dashed lines mark TV distance thresholds of 0.05 and 0.01.}
  \label{fig:im_memory}
\end{figure}

\subsection{Mean time to exit}

\begin{figure}[H]
  \centering
  \includegraphics[width=0.85\textwidth]{memory/mte_10yr.png}
  \caption{Mean time to exit (MTE) by origin state across 10-year birth
           cohorts. MTE is the expected number of transmission steps until
           a lineage's religion first differs from its starting state.}
  \label{fig:mte}
\end{figure}

The individual memory curves (Figure~\ref{fig:im_memory}) translate matrix-level
shifts into a generational time scale. The ``None'' origin state exhibits low
memory in all cohort windows: its TV distance from $\pi^*$ decays rapidly even in
the earliest cohort, reflecting that None-origin lineages begin close to the
implied stationary distribution. The interpretation requires care---low memory for
None-origin lineages indicates proximity to $\pi^*$, which itself assigns
increasing weight to None; it does not necessarily mean frequent exit. For
affiliated states, the curves shift systematically across cohorts toward faster
decay. In the earliest cohorts, Catholic and ``Other'' Protestant origins show the
highest memory; Evangelical and Mainline are roughly paired at intermediate
levels. By the most recent cohort windows, all four affiliated states have
converged toward substantially faster decay, with Catholic declining most sharply.
Mean time to exit (Figure~\ref{fig:mte}) provides a complementary picture.
Catholic MTE falls substantially across cohorts, consistent with the declining
diagonal; Evangelical MTE is comparatively stable, so that the gap between the
two has narrowed considerably by the most recent cohorts. MTE for None-origin
lineages rises across cohorts---an increase commensurate with the rise in None
diagonal persistence from 0.27 to 0.67---a trend that the memory-curve framing
can obscure by conflating persistence with proximity to the stationary
distribution.

\subsection{$\lambda_2$ trend}

\begin{figure}[H]
  \centering
  \includegraphics[width=0.9\textwidth]{memory/lambda2_trend_10yr.png}
  \caption{$\lambda_2$ (modulus of second-largest eigenvalue of $P_t$) by
           10-year birth cohort, with 95\% bootstrap confidence intervals
           (250 replicates, seed 42). A declining $\lambda_2$ indicates that
           intergenerational religious persistence weakens across cohorts.}
  \label{fig:lambda2_trend}
\end{figure}

Figure~\ref{fig:lambda2_trend} presents the headline scalar. $\lambda_2$ declines
approximately linearly from 0.80 in the 1925--34 cohort to 0.60 in the 1975--84
cohort, with 95\% bootstrap confidence intervals. This monotonic decline is the
paper's central empirical finding: intergenerational religious
persistence---measured as the rate at which origin-state information is transmitted
across synthetic generations---has weakened continuously across every successive
cohort window over the six decades covered by the GSS.

\clearpage

%% ═══════════════════════════════════════════════════════════════════════════
\section{Age Effects}
%% ═══════════════════════════════════════════════════════════════════════════

\subsection{Age standardization}

\begin{figure}[H]
  \centering
  \includegraphics[width=0.9\textwidth]{age_std_lambda2_wide}
  \caption{$\lambda_2$ by birth cohort: unstandardized (U), band-restricted
           (R), and age-standardized (S) estimates. Common age band 39--49.
           Convergence of the three series indicates the cohort trend is not
           an age-composition artifact.}
  \label{fig:age_std}
\end{figure}

\begin{figure}[H]
  \centering
  \includegraphics[width=\textwidth]{P_grid_U_vs_S}
  \caption{RELIG16 $\to$ RELIG transition matrices by birth cohort:
           unstandardized (top row) vs.\ age-standardized, band 39--49
           (bottom row). Each column is one cohort window; only cohorts
           whose age support covers the full band are shown.}
  \label{fig:P_grid_U_vs_S}
\end{figure}

A standard concern in cohort analysis is that age-composition differences across
cohort windows could artifactually produce or amplify a trend.
Figures~\ref{fig:age_std} and~\ref{fig:P_grid_U_vs_S} address this directly.
Figure~\ref{fig:age_std} overlays unstandardized (U), band-restricted (R), and
age-standardized (S) $\lambda_2$ series on a common axis, constraining the
comparison to respondents aged 39--49; all three series track closely across
cohorts, and the monotonic decline is preserved in each.
Figure~\ref{fig:P_grid_U_vs_S} presents the underlying transition matrices for
the unstandardized and age-standardized specifications side by side: cell-level
differences are small and do not alter the qualitative pattern. The cohort trend
in $\lambda_2$ is not an age-composition artifact.

\clearpage

%% ═══════════════════════════════════════════════════════════════════════════
\section{Homogeneity Tests}
%% ═══════════════════════════════════════════════════════════════════════════

\subsection{Anderson--Goodman test results}

\begin{table}[htbp]
\centering
\caption{Chi-Square Tests of Homogeneity: Adjacent-Cohort Comparisons}
\label{tab:pairwise}
\begin{tabular}{lrrrrrr}
\toprule
Origin & 1930--40 & 1940--50 & 1950--60 & 1960--70 & 1970--80 & All \\
\midrule
Catholic & 20.8$^{***}$ & 48.2$^{***}$ & 7.2 & 17.5$^{**}$ & 25.2$^{***}$ & 355.3$^{***}$ \\
Evangelical & 20.3$^{***}$ & 27.6$^{***}$ & 14.8$^{**}$ & 12.3$^{*}$ & 23.0$^{***}$ & 313.6$^{***}$ \\
Mainline & 14.3$^{**}$ & 16.0$^{**}$ & 6.6 & 7.9 & 14.4$^{**}$ & 192.9$^{***}$ \\
Other & 3.5 & 11.9$^{*}$ & 5.0 & 6.7 & 10.4$^{*}$ & 123.1$^{***}$ \\
None & 8.2 & 22.0$^{***}$ & 5.6 & 3.4 & 3.3 & 163.8$^{***}$ \\
\midrule
\textbf{Joint} & \textbf{67.0$^{***}$} & \textbf{125.6$^{***}$} & \textbf{39.3$^{**}$} & \textbf{47.8$^{***}$} & \textbf{76.3$^{***}$} & \textbf{1148.7$^{***}$} \\
\bottomrule
\end{tabular}
\begin{minipage}{\linewidth}
\vspace{0.5em}\footnotesize
\textit{Notes:} Cells report $\chi^2$ statistics (Anderson-Goodman 1957).
Each pairwise column tests two consecutive 10-year cohort matrices; df $= (s-1)(T-1) = 4$ per origin (Joint df $= 20$).
The \textit{All} column pools all six cohorts; df $= 20$ per origin (Joint df $= 100$).
Both are reduced where a pooled destination cell is empty.
$^{*}p<0.05$, $^{**}p<0.01$, $^{***}p<0.001$.
\end{minipage}
\end{table}


\subsection{Rolling-window sensitivity}

\begin{figure}[H]
  \centering
  \includegraphics[width=0.85\textwidth]{rolling_window_sensitivity}
  \caption{Sensitivity of the Anderson--Goodman homogeneity test statistic
           to cohort window width. Each point is the chi-square statistic for
           a given window specification; stability across window widths
           confirms the result is not an artifact of binning.}
  \label{fig:rolling}
\end{figure}

The Anderson--Goodman pairwise homogeneity tests (Table~1) formally assess whether
adjacent cohort-specific transition matrices are statistically distinguishable from
one another. \todo{Fill in direction and significance pattern from table before
finalizing.} Figure~\ref{fig:rolling} confirms that these results are not sensitive
to the choice of cohort window width: the chi-square statistics are stable across a
range of binning specifications, ruling out the possibility that heterogeneity is
an artifact of the particular 10-year window used in the main analysis.

\clearpage

%% ═══════════════════════════════════════════════════════════════════════════
\section{Demographic Heterogeneity}
%% ═══════════════════════════════════════════════════════════════════════════

\subsection{Sex}

\begin{figure}[H]
  \centering
  \includegraphics[width=\textwidth]{sex/diff_grid_sex.png}
  \caption{Cell-by-cell difference in transition probabilities (female $-$ male)
           by birth cohort, 1925--1984. Positive values (blue) indicate higher
           conditional probability among female respondents; negative values
           (red) indicate higher probability among male respondents.}
  \label{fig:sex_diff}
\end{figure}

\begin{figure}[H]
  \centering
  \includegraphics[width=\textwidth]{sex/pi_comparison_sex.png}
  \caption{Origin distribution ($\pi_0$) and implied stationary distribution
           ($\pi^*$) by sex and birth cohort, 1925--1984. Bars show
           compositional shares for each religious group within stratum.}
  \label{fig:sex_pi}
\end{figure}

The sex-stratified difference grids (Figure~\ref{fig:sex_diff}) reveal pronounced
heterogeneity by gender. The most salient feature is the strongly male-skewed
inflow into None: in all cohort windows, men exit affiliated traditions into
nonaffiliation at substantially higher rates than women, though this gap attenuates
in the most recent cohorts. In the latest cohort window, the Evangelical column
shows female-skewed inflow---a pattern that may reflect gender-selective processes
of evangelical retention or recruitment in the youngest cohorts. Diagonal
persistence differences between men and women are small and not the primary
interpretive signal. The origin and stationary distribution comparison
(Figure~\ref{fig:sex_pi}) confirms that the male-None and female-Evangelical skews
are reflected in each subgroup's implied long-run distribution: None's $\pi^*$
share is substantially higher for men and Evangelical's $\pi^*$ share higher for
women in recent cohorts.

\subsection{Political views}

\begin{figure}[H]
  \centering
  \includegraphics[width=\textwidth]{politics/diff_grid_polviews_broad.png}
  \caption{Cell-by-cell difference in transition probabilities across political
           views (broad categorization: liberal, moderate, conservative), by
           birth cohort. Each panel shows deviations from the pooled national
           matrix for that cohort.}
  \label{fig:pol_diff}
\end{figure}

\begin{figure}[H]
  \centering
  \includegraphics[width=\textwidth]{politics/pi_comparison_partyid_broad.png}
  \caption{Origin distribution ($\pi_0$) and implied stationary distribution
           ($\pi^*$) by party identification (broad: Democrat, Republican,
           Other) and birth cohort.}
  \label{fig:pol_pi}
\end{figure}

The political heterogeneity results (Figure~\ref{fig:pol_diff}) are consistent
with the broader literature on religious-political sorting: the None column is
disproportionately liberal and the Evangelical column disproportionately
conservative in all cohorts, and both associations sharpen in the most recent
cohort windows. Conservative respondents show somewhat higher diagonal persistence
in the most recent window; this result is worth noting but should not be
overread, given the strong overlap between political conservatism and religious
conservatism in the underlying sample. The party identification comparison
(Figure~\ref{fig:pol_pi}) produces a directionally consistent picture:
Democrat-identified respondents show higher implied None shares in $\pi^*$ and
Republican-identified respondents show higher Evangelical shares, with the gap
widening across cohorts.

\clearpage

%% ═══════════════════════════════════════════════════════════════════════════
\section{Robustness}
%% ═══════════════════════════════════════════════════════════════════════════

\subsection{Non-Black sample}

\begin{figure}[H]
  \centering
  \includegraphics[width=\textwidth]{nonblack/diff_grid_nonblack.png}
  \caption{Cell-by-cell difference in transition probabilities (non-Black $-$
           full sample) by birth cohort. Values near zero indicate that
           Black respondents do not materially shift the national estimates.}
  \label{fig:nb_diff}
\end{figure}

\begin{figure}[H]
  \centering
  \includegraphics[width=\textwidth]{nonblack/im_memory_10yr_nb.png}
  \caption{Individual memory curves (log TV distance from $\pi^*$) for the
           non-Black sample, by origin state and birth cohort.
           Compare to Fig.~\ref{fig:im_memory} for the full-sample baseline.}
  \label{fig:nb_memory}
\end{figure}

The non-Black sample comparison (Figure~\ref{fig:nb_diff}) reveals one substantive
deviation from the full-sample estimates: in the non-Black sample,
Evangelical-to-Mainline outflow is approximately three percentage points higher and
Evangelical diagonal persistence approximately three percentage points lower
throughout the cohort series. This pattern is consistent with the composition of
Black Protestant respondents in the full-sample ``Evangelical'' category and their
distinct transition patterns---a difference partially addressed by the six-state
scheme in the following section. Importantly, the deviation attenuates in later
cohorts, becoming less material by the most recent windows. The memory curves for
the non-Black sample (Figure~\ref{fig:nb_memory}) replicate the national trend
without meaningful deviation: the ordered convergence of affiliated states toward
faster memory decay, and None's persistently low memory throughout, are both
preserved.

\subsection{Six-state scheme (Black Protestant separate)}

\begin{figure}[H]
  \centering
  \includegraphics[width=\textwidth]{bp/P_grid_bp_10yr.png}
  \caption{RELIG16 $\to$ RELIG transition matrices under the 6-state scheme
           (Black Protestant separated from Other Protestant), by 10-year
           birth cohort. Compare to Fig.~\ref{fig:P_grid} for the 5-state
           baseline.}
  \label{fig:bp_grid}
\end{figure}

\begin{figure}[H]
  \centering
  \includegraphics[width=0.85\textwidth]{bp/im_memory_10yr_6state.png}
  \caption{Individual memory curves under the 6-state scheme, by origin state
           and birth cohort. Separation of Black Protestant from Other
           Protestant does not materially alter the national memory trend.}
  \label{fig:bp_memory}
\end{figure}

The six-state scheme, which separates Black Protestant from Other Protestant
(Figure~\ref{fig:bp_grid}), confirms at the cell level the patterns identified in
the non-Black robustness check: Black Protestant and other Evangelical respondents
differ in their outflow profiles, particularly in the Mainline destination. Memory
curves under the six-state scheme (Figure~\ref{fig:bp_memory}) replicate the
national trend in direction and ordering without material change. The headline
result---monotonic decline in intergenerational religious persistence across birth
cohorts---is robust to this state-space extension.

\clearpage

%% ═══════════════════════════════════════════════════════════════════════════
\appendix
\section{Transition Matrix Grids}
%% ═══════════════════════════════════════════════════════════════════════════

\subsection{Individual cohort matrices (5-state, national)}

\begin{figure}[H]
  \centering
  \begin{subfigure}[b]{0.48\textwidth}
    \includegraphics[width=\textwidth]{matrices/trans_1930_10yr.png}
    \caption{Cohort 1925--34}
  \end{subfigure}\hfill
  \begin{subfigure}[b]{0.48\textwidth}
    \includegraphics[width=\textwidth]{matrices/trans_1940_10yr.png}
    \caption{Cohort 1935--44}
  \end{subfigure}
  \begin{subfigure}[b]{0.48\textwidth}
    \includegraphics[width=\textwidth]{matrices/trans_1950_10yr.png}
    \caption{Cohort 1945--54}
  \end{subfigure}\hfill
  \begin{subfigure}[b]{0.48\textwidth}
    \includegraphics[width=\textwidth]{matrices/trans_1960_10yr.png}
    \caption{Cohort 1955--64}
  \end{subfigure}
  \begin{subfigure}[b]{0.48\textwidth}
    \includegraphics[width=\textwidth]{matrices/trans_1970_10yr.png}
    \caption{Cohort 1965--74}
  \end{subfigure}\hfill
  \begin{subfigure}[b]{0.48\textwidth}
    \includegraphics[width=\textwidth]{matrices/trans_1980_10yr.png}
    \caption{Cohort 1975--84}
  \end{subfigure}
  \caption{Individual RELIG16 $\to$ RELIG transition matrices by 10-year
           birth cohort, each panel showing the full $5\times5$ matrix with
           $\pi_0$ and $\pi^*$ bars.}
  \label{fig:app_national_panels}
\end{figure}

\clearpage

\subsection{APC heterogeneity by origin state}

\begin{figure}[H]
  \centering
  \includegraphics[width=0.8\textwidth]{grid_diag_catholic}
  \caption{Catholic diagonal persistence ($P_{\text{cath},\text{cath}}$) by
           birth cohort $\times$ GSS survey period.}
  \label{fig:apc_catholic}
\end{figure}

\begin{figure}[H]
  \centering
  \includegraphics[width=0.8\textwidth]{grid_diag_evangelical}
  \caption{Evangelical diagonal persistence by birth cohort $\times$ survey
           period.}
  \label{fig:apc_evangelical}
\end{figure}

\begin{figure}[H]
  \centering
  \includegraphics[width=0.8\textwidth]{grid_diag_none}
  \caption{Steady-state share of ``None'' ($\pi^*_{\text{none}}$) by birth
           cohort $\times$ survey period.}
  \label{fig:apc_none}
\end{figure}

\begin{figure}[H]
  \centering
  \includegraphics[width=0.8\textwidth]{grid_outflow_to_none}
  \caption{$\pi_0$-weighted outflow to ``None'' from all affiliated origin
           states, by birth cohort $\times$ survey period.}
  \label{fig:apc_outflow}
\end{figure}

\begin{figure}[H]
  \centering
  \includegraphics[width=0.8\textwidth]{grid_lambda2}
  \caption{$\lambda_2$ (second eigenvalue of $P_t$) by birth cohort $\times$
           survey period.}
  \label{fig:apc_lambda2}
\end{figure}

\clearpage

\subsection{Sex-stratified transition matrix grids}

\begin{figure}[H]
  \centering
  \includegraphics[width=\textwidth]{sex/P_grid_male_10yr.png}
  \caption{RELIG16 $\to$ RELIG transition matrices for male respondents, by
           10-year birth cohort, 1925--1984.}
  \label{fig:app_grid_male}
\end{figure}

\begin{figure}[H]
  \centering
  \includegraphics[width=\textwidth]{sex/P_grid_female_10yr.png}
  \caption{RELIG16 $\to$ RELIG transition matrices for female respondents,
           by 10-year birth cohort, 1925--1984.}
  \label{fig:app_grid_female}
\end{figure}

\begin{figure}[H]
  \centering
  \includegraphics[width=\textwidth]{sex/diff_grid_sex.png}
  \caption{Cell-by-cell difference in transition probabilities (female $-$
           male) by birth cohort. Each panel is the element-wise subtraction
           of the male $P_t$ from the female $P_t$.}
  \label{fig:app_sex_diff}
\end{figure}

\begin{figure}[H]
  \centering
  \includegraphics[width=\textwidth]{sex/pi_comparison_sex.png}
  \caption{Origin distribution ($\pi_0$) and implied stationary distribution
           ($\pi^*$) by sex and birth cohort.}
  \label{fig:app_sex_pi}
\end{figure}

\clearpage

\subsection{Political-stratified transition matrix grids}

\begin{figure}[H]
  \centering
  \includegraphics[width=\textwidth]{politics/P_grid_polviews_broad_liberal_10yr.png}
  \caption{RELIG16 $\to$ RELIG transition matrices for liberal respondents
           (polviews broad), by 10-year birth cohort, 1925--1984.}
  \label{fig:app_grid_liberal}
\end{figure}

\begin{figure}[H]
  \centering
  \includegraphics[width=\textwidth]{politics/P_grid_polviews_broad_moderate_10yr.png}
  \caption{RELIG16 $\to$ RELIG transition matrices for moderate respondents
           (polviews broad), by 10-year birth cohort, 1925--1984.}
  \label{fig:app_grid_moderate}
\end{figure}

\begin{figure}[H]
  \centering
  \includegraphics[width=\textwidth]{politics/P_grid_polviews_broad_conservative_10yr.png}
  \caption{RELIG16 $\to$ RELIG transition matrices for conservative respondents
           (polviews broad), by 10-year birth cohort, 1925--1984.}
  \label{fig:app_grid_conservative}
\end{figure}

\begin{figure}[H]
  \centering
  \includegraphics[width=\textwidth]{politics/diff_grid_polviews_broad.png}
  \caption{Cell-by-cell difference in transition probabilities across political
           views (broad), by birth cohort. Each panel shows deviation from the
           pooled national matrix for that cohort.}
  \label{fig:app_pol_diff}
\end{figure}

\clearpage

\subsection{Non-Black sample}

\begin{figure}[H]
  \centering
  \includegraphics[width=\textwidth]{nonblack/P_grid_nonblack_10yr.png}
  \caption{RELIG16 $\to$ RELIG transition matrices for the non-Black sample,
           by 10-year birth cohort, 1925--1984. Compare to
           Fig.~\ref{fig:P_grid} for the full-sample baseline.}
  \label{fig:app_grid_nb}
\end{figure}

\begin{figure}[H]
  \centering
  \includegraphics[width=\textwidth]{nonblack/diff_grid_nonblack.png}
  \caption{Cell-by-cell difference in transition probabilities (non-Black $-$
           full sample) by birth cohort.}
  \label{fig:app_nb_diff}
\end{figure}

\begin{figure}[H]
  \centering
  \includegraphics[width=\textwidth]{nonblack/pi_comparison_nonblack.png}
  \caption{Origin distribution ($\pi_0$) and implied stationary distribution
           ($\pi^*$) for the non-Black subsample vs.\ the full sample, by
           birth cohort.}
  \label{fig:app_nb_pi}
\end{figure}

\clearpage

\subsection{Six-state scheme: full cohort grids}

\begin{figure}[H]
  \centering
  \includegraphics[width=\textwidth]{bp/P_grid_bp_10yr.png}
  \caption{RELIG16 $\to$ RELIG transition matrices under the 6-state scheme
           (Black Protestant as a separate category), by 10-year birth cohort.}
  \label{fig:app_bp_grid}
\end{figure}

\begin{figure}[H]
  \centering
  \includegraphics[width=\textwidth]{bp/pi_dist_grid_bp_10yr.png}
  \caption{Origin distribution ($\pi_0$) and implied stationary distribution
           ($\pi^*$) under the 6-state scheme, by birth cohort.}
  \label{fig:app_bp_pi}
\end{figure}

\subsection{GSS-period stratification}

\begin{figure}[H]
  \centering
  \includegraphics[width=\textwidth]{P_grid_gss_decade}
  \caption{RELIG16 $\to$ RELIG transition matrices by GSS survey period
           (3+2 layout: 1973--1981, 1983--1992, 1993--2002 top row;
           2003--2012, 2013--2020 bottom row). Stratification is by survey
           year of observation rather than birth cohort.}
  \label{fig:app_gss_grid}
\end{figure}

\begin{figure}[H]
  \centering
  \includegraphics[width=\textwidth]{pi_dist_gss_decade}
  \caption{Origin distribution ($\pi_0$, top row) and implied stationary
           distribution ($\pi^*$, bottom row) by GSS survey period. Each bar
           shows the share of respondents in that religious group within the
           given period.}
  \label{fig:app_gss_pi}
\end{figure}

\begin{figure}[H]
  \centering
  \includegraphics[width=0.85\textwidth]{im_memory_gss_decade}
  \caption{Individual memory curves stratified by GSS survey decade (in place
           of birth cohort).}
  \label{fig:app_gss_memory}
\end{figure}

\end{document}
